\documentclass[11pt,a4paper]{article}
\usepackage{fontspec,amsmath,amssymb,amsthm,microtype}
\setmainfont{lmroman10-regular.otf}[BoldFont=lmroman10-bold.otf,ItalicFont=lmroman10-italic.otf,BoldItalicFont=lmroman10-bolditalic.otf]
\usepackage[margin=25mm]{geometry}
\usepackage[unicode,hidelinks]{hyperref}
\setlength{\emergencystretch}{3em}
\setlength{\parskip}{0.4em}
\setlength{\parindent}{0pt}
\allowdisplaybreaks[2]
\raggedbottom
\title{\texorpdfstring{Actual independent-\allowbreak{}source kinetic laws and their Wasserstein infimum}{Actual independent-source kinetic laws and their Wasserstein infimum}}
\author{Omer Deniz Akyildiz\\[0.4em]\normalsize Prepared with AI assistance}
\date{}
\begin{document}\maketitle
Draft mathematical exposition; these notes may contain mistakes.


The position law is an explicit pushforward of an independent product of the initial position law and a fixed source law for common initial velocity and continuous forcing. Every initial coupling is extended by that same source. The proof computes the two output marginals, preserves the initial cost, derives a coupling estimate and takes the infimum. The kinetic pointwise contraction is separately proved from the integral equation and the potential hypotheses; it is not an assumed convergence field.



The proof works over every initial transport plan extended by an independent common source. Exact output marginals and cost control allow the plan cost to be replaced by the true Wasserstein infimum.


\section{\texorpdfstring{Define the independent random-\allowbreak{}map law and compute both output marginals}{Define the independent random-map law and compute both output marginals}}


For a measurable endpoint \(F(x,s)\), set \(\mathcal T_Q\mu=F_\#(\mu\otimes Q)\). Given any coupling \(c\) of \(\mu,\nu\), extend it to \(c\otimes Q\). Projecting to \((y,s)\) produces \(\nu\otimes Q\), because the second marginal of \(c\) is \(\nu\). Pushing further by \(F\) therefore gives \(\mathcal T_Q\nu\). One proves this with the product-map identity, the coupling's marginal equality, and pushforward composition. The corresponding first marginal uses the existing the marginal identity for the independent coupling extension. Both output-law identities are derived from the concrete product construction and are independent of which coupling was chosen.



Given an initial position coupling, add one independent common source containing velocity and forcing. Evaluating the same endpoint map at the two coupled positions gives a coupling of the two desired output laws. Product integration verifies each marginal without assuming the initial coupling is optimal.


\section{\texorpdfstring{Lift pathwise contraction through L\ensuremath{{}^{2}} and take all initial couplings}{Lift pathwise contraction through L² and take all initial couplings}}


Suppose the endpoint obeys \(\|F(x,s)-F(y,s)\|\le q\|x-y\|\) for each common source. On \(c\otimes Q\), the two outputs form a measurable coupling of \(\mathcal T_Q\mu,\mathcal T_Q\nu\). Pointwise contraction passes to the extended L\ensuremath{{}^{2}} norm. The initial difference depends only on \(c\), so its norm on the product is exactly the original coupling cost. Thus \(W_2(\mathcal T_Q\mu,\mathcal T_Q\nu)\le q\,\operatorname{cost}(c)\) for every \(c\). The argument then applies the infimum theorem, using the independent product \(\mu\otimes\nu\) as a witness that couplings exist. No optimal-coupling attainment and no supplied output-law independence are required.



Square the pathwise Lipschitz inequality and integrate over the extended coupling. Taking square roots bounds the output transport cost by the same coefficient times the initial cost. Since the construction works for every initial coupling, taking the infimum gives the Wasserstein inequality.


\section{\texorpdfstring{Retain probability and justify the real-\allowbreak{}valued distance}{Retain probability and justify the real-valued distance}}


A measurable pushforward of \(\mu\otimes Q\) has total mass one when both inputs are probability measures; the probability property of the random-map pushforward records this. The primary distance bound lives in the extended nonnegative reals, so infinite distances are preserved. If the initial Wasserstein distance is finite, the finite coefficient \(q\) makes its upper bound finite, and monotonicity of the finite real value gives the real-valued estimate. One proves that finiteness before converting the inequality, then rewrites the real value of the product. This avoids accidentally converting an infinite output distance to zero.



A measurable map sends a probability measure to a probability measure. Finiteness of the second-moment transport cost is separately before interpreting the extended transport distance as a real number. This prevents real-valued inequalities from concealing an infinite-distance case.


\section{\texorpdfstring{Make the concrete kinetic endpoint a measurable random map}{Make the concrete kinetic endpoint a measurable random map}}


The compact forcing space receives its Borel measurable structure. The source \(s=(v,B)\) contains initial velocity and a continuous forcing path; the separate input is initial position \(x\). the position endpoint map evaluates the selected flow at phase \((v,x)\) and takes its position coordinate. The coordinate rearrangement \((x,(v,B))\mapsto((v,x),B)\) is continuous. Composing it with the proved joint endpoint continuity and projecting position gives continuity, hence measurability, of the map. Independence of source and position is defined later by the explicit product measure rather than encoded as an unspecified output-law premise.



Continuity of the solution with respect to initialization and forcing implies measurability of its endpoint. Thus the law being compared is the pushforward of an explicitly defined random map. Its probability property follows from the input probability measure, not from a formal density calculation.


\section{\texorpdfstring{Derive pointwise contraction from integral solutions}{Derive pointwise contraction from integral solutions}}


Take two values of the selected kinetic flow with phases \((v,x)\) and \((v,y)\) and the same forcing path. Their specification gives the integral equations and equal initial velocities. The proved integral-synchronous theorem cancels their common forcing, derives the right-derivative difference equation, and applies the nonlinear potential energy estimate. Under differentiability, \(m\)-strong convexity, \(M\)-Lipschitz gradient and \(\gamma^2\ge m+M\), this yields \(\|F_t(x,s)-F_t(y,s)\|\le\sqrt2e^{-mt/\gamma}\|x-y\|\). The argument obtains the integral premises by rewriting the flow's initial-value specification, then simplifies the initial position and velocity projections. No differential equation or contraction conclusion is assumed of the chosen flow.



Both endpoints are evaluated with a common initial velocity and common additive path. Their difference satisfies the already proved deterministic kinetic equation. This ties the pointwise Lipschitz coefficient to curvature of the potential, rather than postulating an abstract contractive transition.


\section{\texorpdfstring{Conclude the law-\allowbreak{}to-\allowbreak{}law kinetic theorem}{Conclude the law-to-law kinetic theorem}}


the kinetic position law is the explicit independent-product pushforward for this concrete endpoint. Its probability theorem follows from endpoint measurability. Substituting the derived coefficient \(q=\sqrt2e^{-mt/\gamma}\) in the general random-map transport theorem gives

\[W_2(\mu_t,\nu_t)\le\sqrt2e^{-mt/\gamma}W_2(\mu_0,\nu_0).\]

The source law may be any fixed probability law of initial velocity and continuous forcing, independent of initial position. The statement therefore applies to any such stochastic realization; a stationary Gibbs law is unnecessary for comparing two evolving laws. The extended-valued theorem retains infinite initial costs, and the real-valued theorem explicitly requires finite initial Wasserstein distance. The full all-friction spectral prefactor and a canonical Brownian-source identification are separate refinements beyond this simple-friction law theorem.



The conclusion concerns the independently defined output laws for each initial distribution. Dependence on the auxiliary initial coupling has disappeared through the marginal identities. The coupling argument therefore proves a property of the kinetic transition family.


\section{\texorpdfstring{The full shifted prefactor and exponent pass to output laws}{The full shifted prefactor and exponent pass to output laws}}


The actual-integral general contraction theorem uses the centered gradient sector estimate, so it requires differentiability, strong convexity and a Lipschitz gradient without adding Hessian continuity. For \(d=\gamma-2a>0\), its coefficient is

\[q_a(t)=\frac{\sqrt{2((\gamma-a)^2+a^2)}}{\gamma-2a}
\exp\!\left(\frac{\max\{a^2-m,M-(\gamma-a)^2\}}{\gamma-2a}t\right).\]

The proof instantiates that actual-integral theorem with the same two selected flows and their proved specifications. The coefficient is nonnegative, so it can be bundled as an \(\mathbb R_{\ge0}\) scalar in the already proved random-map transport theorem. This yields \(W_2(\mu_t,\nu_t)\le q_a(t)W_2(\mu_0,\nu_0)\), with the coupling infimum and internally derived output marginals. The real-valued version uses the same finite-initial-cost conversion. The theorem holds for every independent continuous-forcing source law; tying a constructed Brownian source and its required independence to the paper’s named diffusion remains separate, and no Gibbs invariant law is assumed.



The positivity of the transformed-coordinate denominator makes the full prefactor nonnegative. That coefficient can therefore be pulled through the infimum over coupling costs. This yields the entire shifted bound, including its conditioning factor, rather than only the simple zero-shift specialization.


\end{document}
